\subsection{From a marked Kirby diagram to colored ribbon graphs}
\label{subsec:graphical-construction}

\begin{definition}[Twisted categorical input]
\label{def:twisted-categorical-input}
A twisted categorical input is a pair $(\cC,\eta)$ consisting of a
premodular fusion category $\cC$ and a character
\[
  \eta\in A_{\cC}
  =\operatorname{Aut}_{\otimes}(\operatorname{Id}_{\cC}).
\]
For every simple $i$, write
$(\eta)_i=\eta_i\operatorname{id}_i$ as in
\cref{def:character}.  Thus $\eta_i\in\mathbb C^\times$ is a scalar, while
the coupon on an $i$-coloured ribbon is the endomorphism
$\eta_i^n\operatorname{id}_i=(\eta^n)_i$.
\end{definition}

\begin{definition}[Kirby colour]
\label{def:kirby-colour}
The Kirby colour of $\cC$ is the formal linear combination
\[
  \Omega_{\cC}:=
    \sum_{i\in\operatorname{Irr}(\cC)}d_i\,i.
\]
Colouring a link component by $\Omega_{\cC}$ means expanding it into the
displayed finite weighted sum of simple-coloured components.
\end{definition}

\begin{definition}[Tensor-unit projector]
\label{def:tensor-unit-projector}
Let $W$ be an object of $\cC$.  Choose a basis
$u_\mu\colon\mathbf1\to W$ and its composition-dual basis
$\widehat u_\mu\colon W\to\mathbf1$, characterized by
$\widehat u_\mu u_\nu=\delta_{\mu\nu}\operatorname{id}_{\mathbf1}$.
Because $d_{\mathbf1}=1$, this is also the spherical trace-dual basis.  The
projector onto the full
tensor-unit isotypic summand is
\[
  P_{\mathbf1,W}:=\sum_\mu u_\mu\widehat u_\mu
  \in\operatorname{End}_{\cC}(W).
\]
In the dotted construction we write
\[
  \operatorname{Cut}_{\cC}(W):=P_{\mathbf1,W}.
\]
No global-dimension factor is included in this operator.
\end{definition}

\begin{remark}[Basis independence]
\label{rem:tensor-unit-projector-basis-independence}
The endomorphism $P_{\mathbf1,W}$ is independent of the choice of dual
bases.
\end{remark}

\begin{definition}[Exterior closure context]
\label{def:exterior-closure-context}
Let a local ribbon box have input object $U$ and output object $V$.  Fix all
of the coloured tangle outside that box.  After using the standard rigid
identifications at its boundary, the exterior is a morphism
\[
  q_E\colon V\longrightarrow U.
\]
It defines the linear functional
\[
  \operatorname{Cl}_E\colon
  \operatorname{Hom}_{\cC}(U,V)\longrightarrow\mathbb C,
  \qquad
  \operatorname{Cl}_E(f):=\operatorname{tr}_U(q_Ef).
\]
\end{definition}

\begin{remark}[Meaning of exterior closure]
\label{rem:exterior-closure-meaning}
This is exactly the RT evaluation obtained by gluing $f$ into the fixed
exterior tangle.  In particular, an equality $f=g$ in
$\operatorname{Hom}_{\cC}(U,V)$ remains an equality after closure.  We
suppress associators and unitors in this notation by ribbon coherence.
\end{remark}

\begin{definition}[Intrinsic dotted-circle operation]
\label{def:dotted-circle-translation}
Let $\delta$ be a dotted circle in the dotted sublink
$D_{\mathrm{dot}}$ of a special framed link.  In the Crane--Yetter
graphical expansion, $\delta$ means, by definition, the index-$3$ handle
operation of \cite[Definition~5.11, $k=3$]{Tham2021}.  On the boundary-state
object $B_\delta$ of that handle, after resolving the incident ribbons into
simple channels, this operation is
\[
  \operatorname{Cut}_{\cC}(B_\delta)=P_{\mathbf1,B_\delta}.
\]

A dotted co-core chart from \cref{def:dotted-cut-system} is only a way to
display this intrinsic map in a ribbon diagram.  Fix an ordinary-label state
and read the positive collar face as $(p_1,\ldots,p_r)$.  Let $X_j$ be the
label already carried by the oriented ribbon through $p_j$, and put
\[
  X^{[+1]}:=X,\qquad X^{[-1]}:=X^*,
\]
\[
  w_\delta:=((X_1,\sigma_1),\ldots,(X_r,\sigma_r)),
  \qquad
  B_\delta:=|w_\delta|
  :=X_1^{[\sigma_1]}\otimes\cdots\otimes X_r^{[\sigma_r]}.
\]
On the opposite collar face the word is reversed and rigidly dualized, so
canonically
\[
  |w_\delta^{\mathrm{op}}|\cong |w_\delta|^*.
\]
The disk therefore does not change the underlying component label $X_j$;
its sign only specifies how that fixed oriented ribbon occurs in the typed
boundary object.

After all data outside the cut section are fixed, regluing the exterior gives
an exterior closure context in the sense of
\cref{def:exterior-closure-context},
\[
  \mathcal E_\delta\colon
  \operatorname{End}_{\cC}(B_\delta)\longrightarrow\mathbb C.
\]
The dotted contribution is
\[
  \mathcal E_\delta(P_{\mathbf1,B_\delta}).
  \tag{3.1}\label{eq:intrinsic-dotted-context}
\]
After every dotted circle has been translated in this way, the result is a
finite linear combination of genuine closed $\cC$-coloured ribbon graphs in
$S^3$.
\end{definition}

\begin{remark}[Scope of the dotted-circle operation]
\label{rem:dotted-circle-operation-scope}
The dotted circle acts by projection onto the tensor-unit channel.  It is not
itself a coloured ribbon component and is never assigned the ordinary Kirby
colour.  In a modular category the same operation can alternatively be drawn
by encircling with $\operatorname{Dim}(\cC)^{-1}\Omega_{\cC}$, but that
modular replacement is not used here; compare
\cite[Definition~5.11 and Lemma~2.31]{Tham2021}.

The projector is therefore applied to the actual through-channel crossing the
co-core, before the exterior is closed.  It is not postcomposed with a
morphism out of $\mathbf1$: such a postcomposition would be vacuous.  In
particular, if the bunched through-channel is a non-unit simple object $i$,
then $P_{\mathbf1,i}=0$, exactly as required by the index-$3$ operation.
Each resulting closed graph is evaluated by
\cref{def:rt-graph-evaluation}.
\end{remark}

\begin{definition}[Admissible diagrammatic input]
\label{def:admissible-diagrammatic-input}
An admissible input for the Crane--Yetter evaluator consists of
\begin{enumerate}[label=(\roman*),ref=\roman*]
  \item\label{item:evaluator-category}
  a twisted categorical input $(\cC,\eta)$;

  \item\label{item:evaluator-diagram}
  a marked special Kirby representative $(D,\xi_D)$.  A dotted co-core chart may be
  chosen to draw the intrinsic dotted operations, but it is not evaluator
  input;

  \item\label{item:evaluator-cycle}
  a temporary integral cellular $2$-cocycle
  $a=(a_k)\in Z^2(D)\subset C^2(D)$ representing $\xi_D$, where
  $a_k\in\mathbb Z$ is its value on the $k$th primal $2$-handle, represented
  by the $k$th ordinary framed component.
\end{enumerate}
\end{definition}

\begin{remark}[Auxiliary cocycle]
\label{rem:admissible-input-auxiliary-cocycle}
The cocycle in \cref{item:evaluator-cycle} is auxiliary and will be
removed by representative independence.
\end{remark}

\begin{definition}[Role-aware graphical expansion]
\label{def:role-aware-graphical-expansion}
Fix the input of \cref{def:admissible-diagrammatic-input}.  First compose the
intrinsic handle maps belonging to the actual handle decomposition.  A dotted
co-core chart may then be chosen to display those maps.  The resulting
graphical expansion is constructed as follows.
\begin{enumerate}[label=(\roman*),ref=\roman*]
  \item\label{item:ordinary-kirby-expansion}
  Every ordinary framed component is coloured by $\Omega_{\cC}$.  Thus a
  state $\lambda$ chooses a simple label $\lambda_k$ for each such component
  and has the external scalar weight $\prod_k d_{\lambda_k}$.

  \item\label{item:character-coupons}
  On the $k$th ordinary component insert one coupon carrying
  the actual endomorphism
  $(\eta^{a_k})_{\lambda_k}
   =\eta_{\lambda_k}^{a_k}\operatorname{id}_{\lambda_k}$.

  \item\label{item:dotted-cut-projector}
  Apply the intrinsic dotted-circle operation of
  \cref{def:dotted-circle-translation} to every dotted circle.  In a co-core
  chart these disjoint maps are drawn simultaneously.

  \item\label{item:graphical-regluing}
  Expand the projectors using dual bases, interpret all crossings, extrema,
  twists, and coupons by the ribbon graphical calculus, and reglue the cut
  tangle.  Each choice of ordinary labels gives a finite linear combination
  of closed $\cC$-coloured ribbon graphs in $S^3$; its coefficients include
  the fusion-basis coefficients from the tensor-unit projectors.
\end{enumerate}
For a fixed ordinary-label state $\lambda$, let
\[
  \Gamma(D,a;\lambda).
\]
denote the class of this finite linear combination in the complex-linearized
ribbon graphical category.  Graphical evaluation is extended linearly to
$\Gamma(D,a;\lambda)$.  By convention,
$\Gamma(D,a;\lambda)$ contains the projector/fusion expansion coefficients
and all coupon morphisms, but it does \emph{not} contain the external
Kirby-colour coefficient $\prod_kd_{\lambda_k}$; that factor occurs exactly
once in \cref{def:role-aware-handle-evaluator}.
\end{definition}

\begin{remark}[Intrinsic meaning of the graphical expansion]
\label{rem:role-aware-expansion-intrinsic-meaning}
The expression $\Gamma(D,a;\lambda)$ is the intrinsic composite modulo the
ribbon graphical relations, not a chosen planar representative or a chosen
spanning-disk system.  The orientation of a ribbon, not its direction through
a dotted disk, determines whether its label is $X$ or $X^*$.
\end{remark}

\begin{definition}[Homologically twisted Crane--Yetter evaluator]
\label{def:role-aware-handle-evaluator}
For the chosen cocycle $a$, define
\[
  \CY_{\cC,\eta}(D,a)
  :=
  \operatorname{Dim}(\cC)^{1-h_3(D)}
  \sum_{\lambda}
    \left(\prod_k d_{\lambda_k}\right)
    \operatorname{RT}_{\cC}
      \bigl(S^3,\Gamma(D,a;\lambda)\bigr).
\]
\end{definition}

\begin{remark}[Evaluator conventions]
\label{rem:handle-evaluator-conventions}
Here the character coupons are part of the graph
$\Gamma(D,a;\lambda)$; equivalently their simple-component scalars
$\eta_{\lambda_k}^{a_k}$ may be pulled outside the graphical evaluation.
The power $1-h_3(D)$ uses primal handle names and the turned-over convention
of \cref{conv:primal-turned-over}.
\end{remark}

\begin{remark}[A finite sum, not one graph]
\label{rem:evaluator-finite-sum}
The construction does not canonically produce one coloured graph.  It
produces a finite weighted family of graphs because the Kirby colour and the
tensor-unit projectors must be expanded.  Each summand is an honest closed
$\cC$-coloured ribbon graph and is evaluated by
\cref{def:rt-graph-evaluation}.
\end{remark}

\begin{remark}[Proof obligations]
\label{rem:evaluator-proof-obligations}
The preceding definitions do not assume that the scalar is independent of
the cocycle representative, dual bases, graphical presentation of an
intrinsic dotted map, Kirby representative, or Kirby-move sequence.  These
are theorem obligations.  In particular, no evaluation equality is included
in the input data.
\end{remark}
